\chapter{Introduction and Analysis Review}

CS 341 is about \emph{designing} algorithms and \emph{proving} things about them:
that they are correct, and how much time (and occasionally space) they use.
The course is organized around design paradigms (divide and conquer, greedy,
dynamic programming, graph search) and then turns to the limits of efficient
computation (NP-completeness, undecidability). The main reference is CLRS
(\emph{Introduction to Algorithms}, Cormen, Leiserson, Rivest, Stein), with
Kleinberg--Tardos used for some greedy and reduction material.

\section{Model of computation}

To talk about running time we need a model. The default is the \textbf{word RAM}:
\begin{itemize}
  \item memory is an array of words, each holding an integer of $\bigO(\log n)$ bits
        (enough to store an index into the input);
  \item arithmetic ($+,-,\times,\div$), comparisons, and memory reads/writes on single
        words each cost $\bigO(1)$.
\end{itemize}
The running time $T(n)$ of an algorithm is the maximum number of such elementary
operations over all inputs of \emph{size} $n$ (worst-case analysis).

\begin{remark}
Input \emph{size} matters. For a list of $n$ integers the size is usually taken
as $n$ (assuming each fits in a word). For a number-theoretic problem whose input is a
single integer $N$, the size is the number of bits, $\log N$. An algorithm that
runs in time $\bigO(N)$ on input $N$ is then \emph{exponential} in the input size.
This distinction comes back for knapsack (Chapter 5) and subset sum (Chapter 10).
\end{remark}

We also use the \textbf{bit complexity} model occasionally (e.g.\ for integer
multiplication in Chapter 3), where each bit operation costs 1.

\section{Asymptotic notation}

\begin{definition}
Let $f,g:\N\to\R_{\ge 0}$.
\begin{itemize}
  \item $f(n)\in\bigO(g(n))$ if there exist $c>0$ and $n_0$ such that
        $f(n)\le c\,g(n)$ for all $n\ge n_0$.
  \item $f(n)\in\Omega(g(n))$ if there exist $c>0$ and $n_0$ such that
        $f(n)\ge c\,g(n)$ for all $n\ge n_0$.
  \item $f(n)\in\Theta(g(n))$ if $f\in\bigO(g)$ and $f\in\Omega(g)$.
  \item $f(n)\in o(g(n))$ if for \emph{every} $c>0$ there is $n_0$ with
        $f(n)\le c\,g(n)$ for all $n\ge n_0$.
  \item $f(n)\in\omega(g(n))$ if for \emph{every} $c>0$ there is $n_0$ with
        $f(n)\ge c\,g(n)$ for all $n\ge n_0$.
\end{itemize}
\end{definition}

A convenient tool is the \textbf{limit test}: if $L=\lim_{n\to\infty} f(n)/g(n)$ exists, then
\[
L=0 \Rightarrow f\in o(g),\qquad 0<L<\infty \Rightarrow f\in\Theta(g),\qquad L=\infty\Rightarrow f\in\omega(g).
\]
(The converse fails: $f$ can be $\Theta(g)$ even if the limit does not exist,
e.g.\ $f(n)=(2+(-1)^n)n$ and $g(n)=n$.)

\begin{example}
$\log n\in o(n^\varepsilon)$ for every $\varepsilon>0$: by L'H\^opital,
$\lim \frac{\ln n}{n^\varepsilon}=\lim\frac{1/n}{\varepsilon n^{\varepsilon-1}}=\lim\frac{1}{\varepsilon n^\varepsilon}=0$.
Similarly $n^k\in o(c^n)$ for any constant $k$ and $c>1$.
\end{example}

\begin{important}[Growth-rate ladder]
\[
1 \;\prec\; \log\log n \;\prec\; \log n \;\prec\; \sqrt n \;\prec\; n \;\prec\; n\log n \;\prec\; n^2 \;\prec\; n^3 \;\prec\; 2^n \;\prec\; n! \;\prec\; n^n
\]
Also: $\log(n!)\in\Theta(n\log n)$ (Stirling), $\sum_{i=1}^n 1/i = \ln n+\bigO(1)$,
$\sum_{i=1}^n i^k\in\Theta(n^{k+1})$, and a geometric series $\sum_{i=0}^n r^i$ is
$\Theta(1)$ if $r<1$, $\Theta(n)$ if $r=1$, $\Theta(r^n)$ if $r>1$.
\end{important}

Useful properties: $\bigO$ is transitive; $f+g\in\Theta(\max(f,g))$;
$\log_a n\in\Theta(\log_b n)$ for constant bases $a,b>1$, so the base of a log is
irrelevant inside $\bigO$.

\section{Analyzing loops}

For nested loops, write the cost as a sum and evaluate.

\begin{lstlisting}[style=pseudo]
for i = 1 to n do
    for j = i to n do
        for k = i to j do
            $\Theta(1)$ work
\end{lstlisting}
The cost is $\sum_{i=1}^n\sum_{j=i}^n (j-i+1)=\sum_{i=1}^n \binom{n-i+2}{2}\in\Theta(n^3)$.
A quick lower bound argument: for $i\le n/4$ and $j\ge 3n/4$ the inner loop runs at least $n/2$ times,
and there are $(n/4)^2$ such pairs, giving $\Omega(n^3)$.

\begin{aside}
To prove a $\Theta$ bound for a loop, it is often easiest to prove an $\bigO$ upper bound
by bounding every term by the largest, and an $\Omega$ lower bound by keeping only a
constant fraction of the terms, each of which is large.
\end{aside}

\section{The maximum subarray problem}

\begin{definition}[Maximum subarray]
Given an array $A[1..n]$ of (possibly negative) integers, find indices $i\le j$
maximizing $A[i]+A[i+1]+\dots+A[j]$. (We allow the empty subarray with sum $0$, so the answer is never negative.)
\end{definition}

For example, with $A=[2,-6,3,-1,4,-5,2]$ the best is $[3,-1,4]$ with sum $6$.
This problem is the running example for seeing how algorithm design can take us
from $\Theta(n^3)$ to $\Theta(n)$.

\subsection{Attempt 1: brute force, \texorpdfstring{$\Theta(n^3)$}{Theta(n\textasciicircum 3)}}
\begin{lstlisting}[style=pseudo]
best = 0
for i = 1 to n do
    for j = i to n do
        s = 0
        for k = i to j do s = s + A[k]
        best = max(best, s)
return best
\end{lstlisting}
This is exactly the triple loop analyzed above: $\Theta(n^3)$.

\subsection{Attempt 2: reuse partial sums, \texorpdfstring{$\Theta(n^2)$}{Theta(n\textasciicircum 2)}}
The sum $A[i..j]$ equals $A[i..j-1]+A[j]$, so we can maintain it incrementally.
\begin{lstlisting}[style=pseudo]
best = 0
for i = 1 to n do
    s = 0
    for j = i to n do
        s = s + A[j]          // s = A[i] + ... + A[j]
        best = max(best, s)
return best
\end{lstlisting}
Equivalently, precompute prefix sums $P[j]=A[1]+\dots+A[j]$ (with $P[0]=0$); then
the sum of $A[i..j]$ is $P[j]-P[i-1]$.

\subsection{Attempt 3: divide and conquer, \texorpdfstring{$\Theta(n\log n)$}{Theta(n log n)}}
Split $A$ at the middle $m=\floor{(1+n)/2}$. An optimal subarray lies entirely in the left half,
entirely in the right half, or crosses the midpoint. The first two cases are solved
recursively. For the crossing case, the best crossing subarray is
(best suffix of $A[1..m]$) $+$ (best prefix of $A[m+1..n]$), each found by a linear scan.
\begin{lstlisting}[style=pseudo]
MaxSub(A, lo, hi):
    if lo == hi then return max(0, A[lo])
    m = floor((lo + hi) / 2)
    L = MaxSub(A, lo, m);  R = MaxSub(A, m+1, hi)
    s = 0; bestL = 0
    for i = m downto lo do  s = s + A[i]; bestL = max(bestL, s)
    s = 0; bestR = 0
    for j = m+1 to hi do    s = s + A[j]; bestR = max(bestR, s)
    return max(L, R, bestL + bestR)
\end{lstlisting}
Recurrence: $T(n)=2T(n/2)+\Theta(n)$, so $T(n)\in\Theta(n\log n)$ (Chapter 2).

\subsection{Attempt 4: Kadane's algorithm, \texorpdfstring{$\Theta(n)$}{Theta(n)}}
Let $E[j]$ be the maximum sum of a subarray \emph{ending exactly at} $j$ (possibly empty).
Either the best such subarray is empty, or it is $A[j]$ appended to the best subarray ending at $j-1$:
\[
E[0]=0,\qquad E[j]=\max\bigl(0,\;E[j-1]+A[j]\bigr).
\]
The answer is $\max_j E[j]$.
\begin{lstlisting}[style=pseudo]
best = 0; cur = 0
for j = 1 to n do
    cur = max(0, cur + A[j])
    best = max(best, cur)
return best
\end{lstlisting}

\begin{theorem}
Kadane's algorithm returns the maximum subarray sum and runs in $\Theta(n)$ time.
\end{theorem}
\begin{proof}
Runtime is clear: one pass with $\bigO(1)$ work per element.
For correctness, we prove by induction on $j$ the loop invariant: after iteration $j$,
\code{cur} $=E[j]$ and \code{best} $=\max_{0\le t\le j}E[t]$.
Base $j=0$: both are $0=E[0]$. Step: any subarray ending at $j$ is either empty (sum $0$) or
$A[i..j]$ for some $i\le j$; in the latter case its sum is $A[i..j-1]+A[j]\le E[j-1]+A[j]$,
and the value $E[j-1]+A[j]$ is achieved by extending an optimal subarray ending at $j-1$.
Hence $E[j]=\max(0,E[j-1]+A[j])$, which is what \code{cur} becomes. Every subarray ends at some $j$,
so the maximum over all subarrays is $\max_j E[j]$, which is \code{best} at the end.
\end{proof}

Any algorithm must read every element (an unread element could be huge), so $\Omega(n)$
is a lower bound and Kadane is optimal.

\begin{important}[Lesson]
The four solutions illustrate the main themes of the course: a naive algorithm,
then removing redundant work, then divide and conquer, then a dynamic programming
recurrence. The last one is optimal, and proving that needed a loop invariant.
\end{important}

\section{Proving correctness}
Typical tools:
\begin{itemize}
  \item \textbf{Loop invariants}: a statement true before the first iteration (initialization),
        preserved by each iteration (maintenance), and which at termination gives correctness.
  \item \textbf{Induction on input size} for recursive algorithms: assume the recursive calls are
        correct, and show the combine step is correct.
  \item \textbf{Exchange arguments} for greedy algorithms (Chapter 4).
  \item \textbf{Optimal substructure} for dynamic programming (Chapter 5).
\end{itemize}

\section*{Summary}
\begin{itemize}
  \item Running time is measured in the word RAM model, worst case, as a function of input size.
  \item Know the definitions of $\bigO,\Omega,\Theta,o,\omega$, the limit test, and the growth ladder.
  \item Maximum subarray: $\Theta(n^3)$ brute force $\to$ $\Theta(n^2)$ incremental $\to$
        $\Theta(n\log n)$ divide and conquer $\to$ $\Theta(n)$ Kadane.
  \item Correctness is proved with invariants or induction; input size is the number of bits for numeric inputs.
\end{itemize}
